\section{Physical Setup and Hardware Implementation}
\label{sec:ch3-physical-setup}

\subsection{Sensing Subsystems and Hardware Blocks}
\label{sec:ch3-sensing-subsystems}

The physical implementation combined a line-side plug, protected outlet,
solid-state relay module, ESP32-S3 controller, voltage sensor, current sensor,
NTC thermistor, anti-aliasing circuit, backup supply, OLED display, buzzer,
and enclosure. The hardware was organized so that mains conductors, sensing
circuits, low-voltage logic, and user-interface wiring could be traced
separately during testing.

Figure~\ref{fig:ch3-block-diagram} gives the hardware block-level view before
the detailed schematic is introduced.
\begin{figure}[H]
  \centering
  \includegraphics[width=0.98\textwidth]{ch3/figure_3.4.png}
  \caption{Block diagram of the prototype hardware.}
  \label{fig:ch3-block-diagram}
\end{figure}

Figure~\ref{fig:ch3-block-diagram} shows that the relay is positioned as the
final protective actuator, while the voltage, current, and temperature sensors
serve as inputs to the ESP32-S3. This placement allows the firmware to remove
power from the protected outlet after a protection condition is confirmed.

\subsection{Full Schematic Discussion}
\label{sec:ch3-full-schematic-discussion}

The full circuit schematic is shown in Figure~\ref{fig:ch3-full-schematic}.
\clearpage
\begin{figure}[p]
  \centering
  \includegraphics[angle=90,width=0.96\textheight,height=0.96\textwidth,keepaspectratio,trim=0.20in 0.55in 0.20in 0.25in,clip]{ch3/full_schematic.pdf}
  \caption{Full schematic of the TinyML Smart Plug prototype.}
  \label{fig:ch3-full-schematic}
\end{figure}
\clearpage

Figure~\ref{fig:ch3-full-schematic} shows the prototype as several functional
sections rather than as one undifferentiated circuit. The AC input enters the
line-side plug and passes toward the protected outlet through the relay or
load-isolation path, so the relay is the firmware-controlled actuator that
disconnects the load after a latched trip. The voltage-sensing path is taken
from the AC input through the ZMPT101B module and is used only as a monitored
signal for RMS voltage and mains-present logic. The current-sensing path
places the ACS37035 in the load path, routes the conditioned output through
the anti-aliasing and signal-conditioning stage, and digitizes the result with
the MCP3204 ADC for both RMS current and waveform-feature extraction.

The ESP32-S3 section receives the ADC, voltage, and thermistor signals, then
drives the relay-latching circuit, OLED, buzzer, and local indicators. Wi-Fi
communication is handled by the ESP32-S3 and links the local firmware to the
Firebase/PWA monitoring path; the relay decision itself remains local. The
power-supply section separates the low-voltage control electronics from the AC
input through the module supply and local regulators, while signal returns for
the low-voltage circuits are kept within the control side of the prototype.
The schematic supports bench validation of the prototype architecture, but it
does not by itself certify creepage, clearance, enclosure safety, regulatory
compliance, or suitability as a replacement for listed protective devices.
The board-layout views that support this schematic are provided in
Appendix~\ref{app:b}.

\subsection{Signal Conditioning and Enclosure Implementation}
\label{sec:ch3-signal-conditioning-enclosure}

The anti-aliasing circuit and its frequency response are shown in
Figures~\ref{fig:ch3-aaf-circuit} and~\ref{fig:ch3-aaf-bode}.
\begin{figure}[H]
  \centering
  \includegraphics[width=0.70\textwidth]{ch3/figure_3.6.png}
  \caption{Second-order multiple-feedback low-pass anti-aliasing circuit.}
  \label{fig:ch3-aaf-circuit}
\end{figure}

\begin{figure}[H]
  \centering
  \includegraphics[width=0.70\textwidth]{ch3/figure_3.7.png}
  \caption{Bode response of the anti-aliasing filter.}
  \label{fig:ch3-aaf-bode}
\end{figure}

Figure~\ref{fig:ch3-aaf-circuit} shows the analog low-pass filter placed
between the current sensor and the ADC. Figure~\ref{fig:ch3-aaf-bode} verifies
that the filter attenuates higher-frequency content before sampling. The
implemented firmware also applies a digital soft anti-aliasing filter in the
current-feature path, so the analog and digital filters work together rather
than as isolated design elements [57].

The internal arrangement was first modeled before assembly. Figure~\ref{fig:ch3-model-no-case}
shows the prototype layout without the case.
\begin{figure}[H]
  \centering
  \includegraphics[width=0.85\textwidth]{ch3/figure_3.8.png}
  \caption{Three-dimensional model of the prototype without the case.}
  \label{fig:ch3-model-no-case}
\end{figure}

Figure~\ref{fig:ch3-model-no-case} was used to check the relative placement of
the PCB, outlet, relay, fan, and front-panel controls before final enclosure
assembly. This reduced the chance that the enclosure would interfere with
airflow, wiring, or user access.

The corresponding assembled device without the outer case is shown in
Figure~\ref{fig:ch3-actual-no-case}.
\begin{figure}[H]
  \centering
  \includegraphics[width=0.85\textwidth]{ch3/figure_3.9.png}
  \caption{Actual prototype assembly without the case.}
  \label{fig:ch3-actual-no-case}
\end{figure}

Figure~\ref{fig:ch3-actual-no-case} confirms that the physical wiring followed
the planned internal organization. The exposed view was used during bench
testing to probe the sensing, supply, relay, and communication subsystems.

The enclosure model is shown in Figure~\ref{fig:ch3-model-with-case}.
\begin{figure}[H]
  \centering
  \includegraphics[width=0.85\textwidth]{ch3/figure_3.10.png}
  \caption{Three-dimensional model of the prototype with the case.}
  \label{fig:ch3-model-with-case}
\end{figure}

Figure~\ref{fig:ch3-model-with-case} shows the intended user-facing
arrangement: outlet access, front controls, ventilation, and rear plug
connection. The enclosure therefore served both a mechanical and an
operator-interface purpose.

The completed enclosed prototype is shown in Figure~\ref{fig:ch3-actual-with-case}.
\begin{figure}[H]
  \centering
  \includegraphics[width=0.85\textwidth]{ch3/figure_3.11.png}
  \caption{Actual prototype with the case installed.}
  \label{fig:ch3-actual-with-case}
\end{figure}

Figure~\ref{fig:ch3-actual-with-case} shows the plug side and outlet side of
the completed device. This view documents the final physical form used for the
overload, socket-heating, and series arc-fault tests described later in this
chapter.

The major hardware elements are summarized in Table~\ref{tab:ch3-hardware}.
\begin{table}[H]
  \centering
  \caption{Major hardware elements of the prototype}
  \label{tab:ch3-hardware}
  \begin{tabular}{p{0.30\textwidth}p{0.25\textwidth}p{0.35\textwidth}}
    \toprule
    \textbf{Subsystem} & \textbf{Components} & \textbf{Application} \\
    \midrule
    Microcontroller & ESP32-S3 & Local computation, Wi-Fi communication, TinyML inference, relay control \\
    Voltage sensing & ZMPT101B module & RMS voltage measurement and voltage protection input \\
    Current sensing & ACS37035 with MCP3204 ADC & RMS current measurement and waveform-feature source \\
    Temperature sensing & 10k$\Omega$ thermistor & Device temperature measurement for thermal estimation \\
    Protection actuator & Solid-state relay module & Load disconnection after latched protection condition \\
    User interface & OLED, buzzer, LEDs, PWA & Local and remote warning and monitoring \\
    \bottomrule
  \end{tabular}
\end{table}

Table~\ref{tab:ch3-hardware} shows that the prototype combines conventional
sensing and switching hardware with local embedded inference. The hardware
alone does not determine the protection behavior; the firmware computations in
the following sections define how the raw signals are interpreted.
